\documentclass[11pt,a4paper]{article}

% ============================================================
% Paquetes
% ============================================================
\usepackage[utf8]{inputenc}
\usepackage[T1]{fontenc}
\usepackage[spanish]{babel} % <-- CORREGIDO: cambiado a 'spanish'

% --- PAQUETES ---
\usepackage{colortbl}
\usepackage[a4paper,margin=2.5cm]{geometry}
\usepackage{graphicx}
\usepackage{array}
\usepackage{longtable}
\usepackage{tabularx}
\usepackage{booktabs}
\usepackage{xcolor}

% --- SOLUCIÓN DEFINITIVA PARA EL DESBORDAMIENTO DE URL ---
\usepackage{xurl} % Permite saltos de línea en CUALQUIER carácter de la URL

% Carga hyperref DESPUÉS de xurl
\usepackage[colorlinks=true, linkcolor=black, urlcolor=blue, citecolor=black]{hyperref}

\usepackage{enumitem}
\usepackage{titlesec}
\usepackage{fancyhdr}
\usepackage{listings}
\usepackage{tcolorbox}
\tcbuselibrary{skins,breakable}

% --- REDEFINICIONES ---
\renewcommand{\contentsname}{Contenido}
\renewcommand{\tablename}{Tabla}
\renewcommand{\figurename}{Figura}

% ============================================================
% Colores
% ============================================================
\definecolor{primario}{HTML}{2C3E50}
\definecolor{secundario}{HTML}{3498DB}
\definecolor{exito}{HTML}{27AE60}
\definecolor{alerta}{HTML}{E67E22}
\definecolor{peligro}{HTML}{C0392B}
\definecolor{gris}{HTML}{7F8C8D}
\definecolor{grisclaro}{HTML}{ECF0F1}
\definecolor{azulclaro}{HTML}{EBF5FB}
\definecolor{verdeclaro}{HTML}{E8F8F0}
\definecolor{naranjaclaro}{HTML}{FDF2E9}
\definecolor{rojoclaro}{HTML}{FDEDEC}

% ============================================================
% Hyperref
% ============================================================
\hypersetup{
    colorlinks=true,
    linkcolor=secundario,
    urlcolor=secundario,
    pdftitle={Instructivo de Usuario - Sistema EME},
    pdfauthor={Edson Joel Carrera Avila},
    pdfsubject={Instructivo de instalacion y uso}
}

% ============================================================
% Titulos
% ============================================================
\titleformat{\section}
    {\normalfont\Large\bfseries\color{primario}}
    {\thesection.}{0.5em}{}
\titleformat{\subsection}
    {\normalfont\large\bfseries\color{primario}}
    {\thesubsection.}{0.5em}{}
\titleformat{\subsubsection}
    {\normalfont\normalsize\bfseries\color{primario}} % <-- CORREGIDO: Se eliminó el salto de línea que causaba error
    {\thesubsubsection.}{0.5em}{}

% ============================================================
% Encabezado / pie
% ============================================================
\pagestyle{fancy}
\fancyhf{}
\fancyhead[L]{\small\textit{Expediente Médico Electrónico}}
\fancyhead[R]{\small\textit{Instructivo de Usuario}}
\fancyfoot[C]{\thepage}
\renewcommand{\headrulewidth}{0.4pt}

% ============================================================
% Listas
% ============================================================
\setlist[itemize]{leftmargin=*,itemsep=3pt,parsep=0pt,topsep=4pt}
\setlist[enumerate]{leftmargin=*,itemsep=4pt,parsep=0pt,topsep=4pt}

% ============================================================
% Tipos de columna
% ============================================================
\newcolumntype{L}[1]{>{\raggedright\arraybackslash}p{#1}}
\newcolumntype{C}[1]{>{\centering\arraybackslash}p{#1}}

% ============================================================
% Cajas de aviso
% ============================================================
\newtcolorbox{cajaInfo}{
  enhanced, breakable,
  colback=azulclaro, colframe=secundario,
  boxrule=0.5pt, arc=2pt, left=8pt, right=8pt, top=6pt, bottom=6pt,
  before skip=8pt, after skip=8pt
}
\newtcolorbox{cajaTip}{
  enhanced, breakable,
  colback=verdeclaro, colframe=exito,
  boxrule=0.5pt, arc=2pt, left=8pt, right=8pt, top=6pt, bottom=6pt,
  before skip=8pt, after skip=8pt
}
\newtcolorbox{cajaAlerta}{
  enhanced, breakable,
  colback=naranjaclaro, colframe=alerta,
  boxrule=0.5pt, arc=2pt, left=8pt, right=8pt, top=6pt, bottom=6pt,
  before skip=8pt, after skip=8pt
}
\newtcolorbox{cajaCritica}{
  enhanced, breakable,
  colback=rojoclaro, colframe=peligro,
  boxrule=0.8pt, arc=2pt, left=8pt, right=8pt, top=6pt, bottom=6pt,
  before skip=8pt, after skip=8pt
}

% ============================================================
% Listings (codigo / consola)
% ============================================================
\lstset{
    basicstyle=\ttfamily\small,
    breaklines=true,
    columns=fullflexible,
    backgroundcolor=\color{grisclaro},
    frame=single,
    framerule=0.5pt,
    rulecolor=\color{gris}
}

% ============================================================
% Información del documento
% ============================================================
\begin{document}

\begin{titlepage}
    \centering
    \vspace*{1.5cm}

    {\Large \textbf{Universidad Autónoma de la Ciudad de México}}\\[0.3cm]
    {\large Licenciatura en Ingeniería de Software}\\[0.2cm]
    {\large Análisis de Requisitos}\\[2.5cm]

    \rule{\linewidth}{0.5pt}\\[0.4cm]
    {\LARGE \bfseries Segundo Avance de Proyecto}\\[0.3cm]
    {\Large Expediente Médico Electrónico}\\[0.3cm]
    \rule{\linewidth}{0.5pt}\\[2cm]

    \begin{flushleft}
        \textbf{Alumno:} Edson Joel Carrera Avila \\[0.3cm]
        \textbf{Alumno:} Mauricio Bonifacio Hernández \\[0.3cm]
        \textbf{Grupo:} 302 \\[0.3cm]
        \textbf{Profesor:} Prof. Mtro. Raul Rolando Jara Montes \\[0.3cm]
    \end{flushleft}
    \vfill
\end{titlepage}

% <-- CORREGIDO: Aislamiento del color del índice para mantenerlo negro
{
    \hypersetup{linkcolor=black}
    \tableofcontents
}
\newpage

% ============================================================
% 1. QUE ES EL SISTEMA
% ============================================================
\section{Que es este sistema?}

\textbf{EME (Expediente Medico Electronico)} es una aplicacion web disenada para consultorios medicos pequenos. Permite gestionar de forma centralizada y segura toda la informacion clinica y administrativa de los pacientes.
\subsection*{El sistema permite:}

\begin{itemize}
    \item Registrar pacientes y abrirles un expediente clinico unico.
    \item Agendar citas medicas con un flujo centrado en el paciente.
    \item Registrar consultas con signos vitales, diagnostico y recetas.
    \item Manejar tres tipos de usuarios con permisos diferenciados: \textbf{Administrador, Medico y Recepcionista}.
    \item Llevar un registro automatico (auditoria) de todo lo que ocurre en el sistema.
\end{itemize}

\begin{cajaInfo}
    \textbf{Importante:} el sistema funciona \textbf{en tu propia computadora} (en \texttt{localhost}), no en internet. Es ideal para un consultorio pequeno donde la informacion se mantiene local.
\end{cajaInfo}

% ============================================================
% 2. REQUISITOS PREVIOS
% ============================================================
\section{Requisitos previos}

Solo necesitas dos cosas instaladas en tu computadora:

\begin{center}
    \begin{tabular}{|L{3.5cm}|L{3cm}|L{6cm}|}
        \hline
        \rowcolor{grisclaro}
        \textbf{Requisito} & \textbf{Version} & \textbf{Donde se descarga}                          \\
        \hline
        Node.js
                           & v18 o superior   & \url{https://nodejs.org/es/download/}               \\
        \hline
        Navegador web      & Reciente         & Chrome, Firefox, Edge o Safari (ya viene instalado) \\
        \hline
    \end{tabular}
\end{center}

\begin{cajaTip}
    \textbf{npm} se instala automaticamente junto con Node.js, no necesitas instalarlo aparte.
\end{cajaTip}

\begin{enumerate}
    \item Abre la \textbf{Terminal} (Linux/macOS) o \textbf{Simbolo del sistema} (Windows):
          \begin{itemize}
              \item \textbf{Windows:} pulsa la tecla \texttt{Windows + R}, escribe \texttt{cmd} y presiona Enter.
              \item \textbf{Linux:} abre tu Terminal habitual.
          \end{itemize}
    \item Escribe este comando y presiona Enter:
          \begin{lstlisting}
node --version
\end{lstlisting}
    \item Si aparece algo como \texttt{v18.17.0} o un numero mayor, \textbf{ya lo tienes}. Salta al Paso 1.
    \item Si aparece \textit{"comando no reconocido"} o un error, descarga e instala Node.js desde el enlace anterior. \textbf{Cierra y vuelve a abrir la Terminal} despues de instalar.
\end{enumerate}

% ============================================================
% 3. PASO 1
% ============================================================
\section{Descargar y preparar el proyecto}

\begin{enumerate}
    \item Descarga el archivo \url{https://github.com/5to-AnalisisRequisitos/Proyecto01_ExpedienteMedicoElectronico/archive/refs/heads/main.zip}.
    \item El archivo comprimido, sin importar el sistema operativo, se decargara automaticamente en la carpeta "Descargas".
    \item Descomprimelo en la ruta que quieras, debes ver una carpeta con esta estructura:
\end{enumerate}

\begin{lstlisting}
Proyecto01_ExpedienteMedicoElectronico-main/
|-- server.js          
|-- setup.js              
|-- package.json        
|-- README.md            
|-- iniciar.sh           
|-- iniciar.bat                 
|-- .gitignore        
\-- public/
    |-- index.html    
    |-- styles.css         
    \-- app.js          
\end{lstlisting}

\begin{cajaCritica}
    \textbf{No muevas, renombres ni borres ningun archivo.} El sistema necesita todos para funcionar correctamente.
\end{cajaCritica}

% ============================================================
% 4. PASO 2: INICIAR
% ============================================================
\section{Iniciar el sistema}

Hay un script preparado que hace todo el trabajo automaticamente: verifica Node.js, instala lo necesario y arranca el servidor.
\subsection{Usuario de Windows}

\begin{enumerate}
    \item Abre la \textbf{Terminal} dentro de la carpeta del proyecto (clic derecho sobre la carpeta y \textit{"Abrir en terminal"}.
    \item Ejecuta el script:
          \begin{lstlisting}
./iniciar.bat
        \end{lstlisting}
    \item Cuando termine, veras un mensaje parecido a este:
\end{enumerate}

\begin{lstlisting}
========================================
  EME v2.0 - Expediente Medico Electronico
  Conforme a IEEE 830 + Caso de Uso
========================================
  Servidor:    http://localhost:3000
  Estado:      Funcionando
  JWT Secret:  Configurado (.env)
  Rate limit:  5 intentos / 15 min
  Helmet:      Cabeceras de seguridad activas
  CORS:        Restringido a localhost:3000
\end{lstlisting}

\begin{cajaCritica}
    \textbf{No cierres esa ventana negra} mientras estes usando el sistema. Si la cierras, el sistema se apaga.
\end{cajaCritica}

\subsection{Usuario de Linux}

\begin{enumerate}
    \item Abre la \textbf{Terminal} dentro de la carpeta del proyecto (clic derecho sobre la carpeta y \textit{"Abrir terminal aqui"}.
    \item \textbf{Solo la primera vez}, da permisos de ejecucion al script:
          \begin{lstlisting}
chmod +x iniciar.sh
\end{lstlisting}
    \item Ejecuta el script:
          \begin{lstlisting}
./iniciar.sh
\end{lstlisting}
    \item Veras el mismo mensaje de confirmacion que aparece arriba.
\end{enumerate}

\begin{cajaCritica}
    \textbf{No cierres la terminal} mientras uses el sistema.
\end{cajaCritica}

\begin{cajaTip}
    Si aparece un mensaje sobre que falta Node.js, el script te abrira automaticamente la pagina oficial de descarga. Instalalo, cierra y vuelve a abrir todo, y ejecuta de nuevo el script.
\end{cajaTip}

% ============================================================
% 5. PASO 3: NAVEGADOR
% ============================================================
\section{Abrir el sistema en el navegador}

\begin{enumerate}
    \item Abre tu navegador favorito (Chrome, Firefox, Edge o Safari).
    \item En la barra de direcciones escribe \textbf{exactamente}:
          \begin{lstlisting}
http://localhost:3000
\end{lstlisting}
    \item Presiona Enter.
    \item Listo. Ya estas dentro del sistema.
\end{enumerate}

% ============================================================
% 6. PASO 4: PRIMER ADMIN
% ============================================================
\section{Crear el primer Administrador}

\begin{cajaInfo}
    \textbf{Esta seccion solo aplica la primera vez que abres el sistema.} Despues de crear el primer Administrador, esta opcion se desactiva.
\end{cajaInfo}

La primera vez que abres el sistema \textbf{no hay ningun usuario creado}, asi que veras la pantalla \textit{"Bienvenido al EME"} que te invita a configurar el sistema.
\begin{enumerate}
    \item Haz clic en el boton \textbf{"Crear administrador"}.
    \item Llena el formulario con tus datos:
\end{enumerate}

\begin{center}
    \begin{tabular}{|L{3.5cm}|L{9cm}|}
        \hline
        \rowcolor{grisclaro}
        \textbf{Campo}       & \textbf{Que poner}                                                         \\
        \hline
        Nombre               & Tu nombre                                                                  \\
        \hline
        Apellido             & Tu apellido                                                                \\
        \hline
        Email                & Tu correo electronico (sera tu usuario para entrar)                        \\
        \hline
        Telefono             & Tu telefono                                                                \\
        \hline
        Contrasena           & Minimo 8 caracteres, con al menos una mayuscula, una minuscula y un numero \\
        \hline
        Confirmar contrasena & Repite exactamente la misma contrasena                                     \\
        \hline
    \end{tabular}
\end{center}

\begin{enumerate}
    \setcounter{enumi}{2}
    \item Haz clic en \textbf{"Crear administrador"}.
    \item Veras la pantalla \textit{"!Cuenta creada!"} que muestra un \textbf{codigo unico de recuperacion} con esta forma:
          \begin{lstlisting}
EME-XXXX-XXXX-XXXX-XXXX
\end{lstlisting}
\end{enumerate}

\begin{cajaCritica}
    \textbf{IMPORTANTE --- leer con atencion:}
    \begin{itemize}
        \item Este codigo aparece \textbf{una sola vez}, ya no podras verlo otra vez.
        \item Es \textbf{la unica forma} de recuperar tu cuenta si olvidas la contrasena.
        \item Tienes tres opciones para guardarlo, los botones estan en la misma pantalla:
              \begin{itemize}
                  \item \textbf{Copiar} al portapapeles.
                  \item \textbf{Descargar} como archivo de texto.
                  \item Escribirlo en un papel o un gestor de contrasenas.
              \end{itemize}
        \item \textbf{Guardalo en un lugar seguro} (cajon con llave, gestor de contrasenas, archivo cifrado).
    \end{itemize}
\end{cajaCritica}

\begin{enumerate}
    \setcounter{enumi}{4}
    \item Marca la casilla \textbf{"He guardado mi codigo de recuperacion"}.
    \item Haz clic en \textbf{"Continuar al inicio de sesion"}.
    \item Ya en la pantalla de \textbf{Login}, escribe tu correo y contrasena y entra al sistema.
\end{enumerate}

% ============================================================
% 7. PASO 5: CREAR EQUIPO
% ============================================================
\section{Crear las cuentas de tu equipo}

Una vez dentro como Administrador, lo primero que conviene hacer es crear las cuentas para el resto del personal.
\subsection{Los tres roles del sistema}

% <-- CORREGIDO: Se eliminó el envoltorio \begin{center} y \end{center} para evitar las advertencias.
\begin{longtable}{|L{2.8cm}|L{5cm}|L{5cm}|}
    \hline
    \rowcolor{grisclaro}
    \textbf{Rol}           & \textbf{Que puede hacer}                                                                               & \textbf{Que NO puede hacer}                                               \\
    \hline
    \endhead

    \textbf{Administrador} & Crear y gestionar cuentas, ver auditoria global, restaurar pacientes eliminados, ver estadisticas.     & Ver expedientes, consultas ni datos clinicos.                             \\
    \hline

    \textbf{Medico}        & Todo lo relacionado a pacientes: registrar, consultar, atender consultas, recetar, eliminar pacientes. & Crear cuentas de usuario.                                                 \\
    \hline

    \textbf{Recepcionista} & Registrar pacientes nuevos, agendar y cancelar citas, editar datos basicos.                            & Ver consultas, diagnosticos ni recetas. Tampoco puede eliminar pacientes. \\
    \hline
\end{longtable}

\subsection{Como crear un usuario nuevo}

\begin{enumerate}
    \item En el menu superior, entra a \textbf{"Gestion de usuarios"}.
    \item Haz clic en el boton azul \textbf{"+ Crear usuario"}.
    \item Llena el formulario:
          \begin{itemize}
              \item \textbf{Nombre y apellido} del usuario.
              \item \textbf{Email} (sera su usuario para entrar).
              \item \textbf{Rol}: elige Administrador, Medico o Recepcionista.
              \item \textbf{Especialidad y numero de cedula profesional}: solo aparecen si elegiste Medico.
              \item \textbf{Contrasena temporal} (la persona puede cambiarla despues).
          \end{itemize}
    \item Haz clic en \textbf{"Crear"}.
    \item El sistema generara un \textbf{codigo unico de recuperacion} para ese usuario.
\end{enumerate}

\begin{cajaAlerta}
    \textbf{Copia el codigo y entregaselo a la persona en privado.} Es su responsabilidad guardarlo en un lugar seguro. Repite este paso por cada persona del equipo.
\end{cajaAlerta}

% ============================================================
% 8. USO MEDICO
% ============================================================
\section{Uso diario: Medico}

Cuando un medico inicia sesion, ve un \textbf{Dashboard} con sus estadisticas: citas de hoy, total de pacientes, expedientes abiertos y consultas del mes.
\subsection{Registrar un paciente nuevo}

\begin{enumerate}
    \item Entra a \textbf{"Pacientes"} en el menu.
    \item Haz clic en \textbf{"+ Nuevo paciente"}.
    \item Llena los datos: CURP, nombre, apellido, fecha de nacimiento, sexo, telefono, etc.
    \item Escribe el \textbf{motivo de apertura del expediente} (ej. \textit{"Primera consulta"}).
    \item Si quieres, agrega tipo de sangre, alergias y observaciones (seccion colapsable).
    \item Haz clic en \textbf{"Crear paciente y expediente"}.
    \item El sistema generara automaticamente un \textbf{numero de expediente unico} con formato \texttt{EXP-2026-00001}.
\end{enumerate}

\begin{cajaTip}
    Si la CURP ya existe en el sistema, te avisara y mostrara al paciente existente para evitar duplicados.
\end{cajaTip}

\subsection{Buscar un paciente}

En la seccion \textbf{"Pacientes"} usa el campo \textbf{"Buscar por nombre, CURP o telefono..."}. Empieza a escribir y los resultados apareceran automaticamente.

\subsection{Agendar una cita}

\begin{enumerate}
    \item Abre el \textbf{detalle del paciente} haciendo clic sobre su nombre en la tabla.
    \item Haz clic en \textbf{"Agendar cita"}.
    \item Selecciona la \textbf{fecha} (no puede ser anterior a hoy), la \textbf{hora}, el \textbf{medico} y el \textbf{motivo}.
    \item Haz clic en \textbf{"Agendar"}.
\end{enumerate}

\subsection{Registrar una consulta}

\begin{enumerate}
    \item Entra a \textbf{"Citas"} o al detalle del paciente.
    \item Haz clic en \textbf{"Atender"} o \textbf{"Nueva consulta"} en la cita correspondiente.
    \item Llena los \textbf{signos vitales}: peso, altura, presion arterial, temperatura, frecuencia cardiaca.
    \item Escribe el \textbf{motivo de consulta}, \textbf{sintomas}, \textbf{diagnostico} y \textbf{plan de tratamiento}.
    \item Si quieres, agrega \textbf{medicamentos prescritos} con dosis, cantidad y duracion.
    \item Haz clic en \textbf{"Guardar consulta"}.
\end{enumerate}

La consulta queda registrada y la cita se marca automaticamente como \textit{"completada"}.
\subsection{Eliminar un paciente (solo medicos)}

\begin{enumerate}
    \item Abre el detalle del paciente.
    \item Haz clic en \textbf{"Eliminar paciente"}.
    \item Escribe un \textbf{motivo obligatorio} de al menos 5 caracteres.
    \item Confirma.
\end{enumerate}

\begin{cajaInfo}
    El paciente se marca como inactivo (\textbf{no se borra realmente}), su expediente se archiva y sus citas pendientes se cancelan automaticamente. Solo el Administrador podra restaurarlo.
\end{cajaInfo}

% ============================================================
% 9. USO RECEPCIONISTA
% ============================================================
\section{Uso diario: Recepcionista}

La recepcionista usa el sistema principalmente para registrar pacientes y agendar citas.
\begin{itemize}
    \item \textbf{Puede:} registrar pacientes, editar sus datos (con motivo obligatorio), agendar y cancelar citas, ver la agenda del dia.
    \item \textbf{No puede:} ver consultas, diagnosticos, recetas ni antecedentes medicos.
\end{itemize}

Los pasos son los mismos que para el medico en las secciones de \textbf{Registrar paciente}, \textbf{Buscar paciente} y \textbf{Agendar cita} (puntos 8.1 a 8.3 de este instructivo).
% ============================================================
% 10. USO ADMIN
% ============================================================
\section{Uso diario: Administrador}

El administrador no atiende pacientes. Su trabajo es gestionar usuarios y supervisar el sistema.

\subsection{Secciones del menu}

\begin{center}
    \begin{tabular}{|L{5cm}|L{8cm}|}
        \hline
        \rowcolor{grisclaro}
        \textbf{Seccion}             & \textbf{Para que sirve}                                             \\
        \hline
        Dashboard                    & Resumen estadistico del sistema                                     \\
        \hline
        Gestion de usuarios          & Crear, editar, activar o desactivar cuentas                         \\
        \hline
        Auditoria global del sistema & Ver todo lo que pasa: logins, accesos denegados, cambios            \\
        \hline
        Pacientes eliminados         & Ver pacientes con eliminacion logica y restaurarlos si es necesario \\
        \hline
        Mi actividad                 & Ver tus propias acciones registradas                                \\
        \hline
    \end{tabular}
\end{center}

\subsection{Desactivar a un usuario}

\begin{enumerate}
    \item Entra a \textbf{"Gestion de usuarios"}.
    \item Busca al usuario en la tabla.
    \item Haz clic en \textbf{"Desactivar"}.
    \item El usuario no podra iniciar sesion hasta que lo vuelvas a activar.
\end{enumerate}

\begin{cajaTip}
    No puedes desactivarte a ti mismo, ni desactivar al ultimo Administrador activo del sistema.
\end{cajaTip}

\subsection{Restaurar un paciente eliminado}

\begin{enumerate}
    \item Entra a \textbf{"Pacientes eliminados"}.
    \item Busca el paciente.
    \item Haz clic en \textbf{"Restaurar"}.
    \item El paciente vuelve a estar activo, igual que su expediente.
\end{enumerate}

\subsection{Consultar la auditoria global}

\begin{enumerate}
    \item Entra a \textbf{"Auditoria global del sistema"}.
    \item Puedes filtrar por:
          \begin{itemize}
              \item \textbf{Tipo de accion} (Login, Login fallido, Crear, Actualizar, Eliminar, etc.).
              \item \textbf{Email del usuario} (escribe parte del email, ej. \texttt{dr@}).
          \end{itemize}
    \item Haz clic en \textbf{"Filtrar"}.
\end{enumerate}

% ============================================================
% 11. RECUPERAR CONTRASENA
% ============================================================
\section{Recuperar contrasena olvidada}

Si olvidas tu contrasena y tienes tu codigo de recuperacion:

\begin{enumerate}
    \item En la pantalla de Login, haz clic en \textbf{"!Olvidaste tu contrasena?"}.
    \item Escribe tu \textbf{email} y tu \textbf{codigo de recuperacion} (\texttt{EME-XXXX-XXXX-XXXX-XXXX}).
    \item Haz clic en \textbf{"Verificar codigo"}.
    \item Te llevara a la pantalla \textbf{"Nueva contrasena"}.
    \item Escribe la contrasena nueva dos veces (con los mismos requisitos: minimo 8 caracteres, mayuscula, minuscula y numero).
    \item Haz clic en \textbf{"Restablecer contrasena"}.
\end{enumerate}

\begin{cajaCritica}
    \textbf{El sistema te dara un nuevo codigo de recuperacion.} Guardalo igual que la primera vez. El anterior queda inservible.
\end{cajaCritica}

\begin{cajaInfo}
    Si \textbf{perdiste el codigo de recuperacion}, pide al Administrador que reasigne tu cuenta. Si quien perdio el codigo es el unico Administrador, sera necesario reiniciar el sistema (ver Preguntas Frecuentes).
\end{cajaInfo}

% ============================================================
% 12. APAGAR
% ============================================================
\section{Como apagar el sistema}

\begin{enumerate}
    \item Asegurate de que todos los usuarios hayan cerrado sesion.
    \item Ve a la ventana negra (la consola que se abrio al iniciar).
    \item Presiona \texttt{Ctrl + C} en tu teclado.
    \item Confirma con \texttt{S} o \texttt{Y} si lo pregunta.
    \item Cierra la ventana.
\end{enumerate}

\begin{cajaTip}
    Mientras la ventana negra este abierta, el sistema sigue funcionando. Si solo cierras el navegador, el servidor sigue activo.
\end{cajaTip}

% ============================================================
% 13. FAQ
% ============================================================
\section{Preguntas frecuentes y solucion de problemas}

\subsection*{Al abrir \texttt{http://localhost:3000} el navegador dice \textit{"No se puede acceder"}}

Significa que el servidor no esta corriendo. Vuelve al \textbf{Paso 2} y ejecuta \texttt{iniciar.bat} o \texttt{iniciar.sh} otra vez. Asegurate de no cerrar la ventana negra.
\subsection*{Aparece \textit{"Demasiados intentos. Espera 15 minutos"}}

Has fallado 5 veces seguidas al iniciar sesion. El sistema te bloquea durante 15 minutos para evitar ataques. Espera y vuelve a intentar; revisa que el email y la contrasena sean correctos.

\subsection*{La sesion se cerro sola}

El sistema \textbf{cierra automaticamente la sesion tras 15 minutos sin actividad}. Antes de cerrar veras un aviso con cuenta regresiva de 60 segundos: si pulsas \textbf{"Seguir conectado"}, la sesion se mantiene. Esta es una medida de seguridad para que nadie use tu cuenta si te alejas de la computadora.
\subsection*{Olvide el codigo de recuperacion del unico Administrador}

En este caso, no hay forma de recuperar la cuenta. Tendras que reiniciar el sistema:

\begin{enumerate}
    \item Apaga el servidor (\texttt{Ctrl + C} en la consola).
    \item \textbf{Haz una copia de respaldo} del archivo \texttt{eme.db} por si necesitas consultar los datos antiguos.
    \item Borra el archivo \texttt{eme.db} de la carpeta del proyecto.
    \item Vuelve a iniciar el sistema con \texttt{iniciar.bat} o \texttt{iniciar.sh}.
    \item El sistema te pedira crear un nuevo Administrador. Esta vez guarda muy bien el codigo.
\end{enumerate}

\begin{cajaCritica}
    Esto \textbf{borrara todos los datos del sistema} (pacientes, citas, consultas). Solo es necesario en ultimo caso.
\end{cajaCritica}

\subsection*{Donde se guardan los datos?}

Todos los datos estan en el archivo \texttt{eme.db} dentro de la carpeta del proyecto. Es una base de datos SQLite. Si quieres hacer respaldos, copia ese archivo a un lugar seguro periodicamente.
\subsection*{Puedo usar el sistema desde otra computadora en la misma red?}

El sistema fue disenado para funcionar en \textbf{localhost} (solo la computadora donde corre el servidor). Por seguridad (CORS restringido), no se permite acceso desde otras computadoras de la red sin modificaciones adicionales en \texttt{server.js}.
\subsection*{Puedo cambiar el puerto 3000 por otro?}

Si. Abre el archivo \texttt{.env} (se genera automaticamente la primera vez), y cambia el valor de \texttt{PORT=3000} por el puerto que quieras (por ejemplo \texttt{PORT=8080}). Luego reinicia el sistema y accede desde \texttt{http://localhost:8080}.

\subsection*{La instalacion se queda atorada en \texttt{npm install}}

Verifica tu conexion a internet (se descargan los paquetes de Node.js desde \texttt{npmjs.com}). Si tienes proxy o firewall corporativo, podria estar bloqueando la descarga. Pide ayuda a tu administrador de TI.

\end{document}